\chapter{Introduction}\label{ch:introduction}

\section{Motivation}\label{sec:intro:motivation}

A robot that plans with a generative model gains expressiveness at a cost in computation. Diffusion
planners \parencite{janner2022planning,chi2023diffusion} represent multi-modal demonstration data
\parencite{jia2024towards} well and, when paired with a constrained projection step
\parencite{romer2025diffusion}, can respect dynamics and safety at deployment. The price is the sampling
loop: producing one plan costs tens to hundreds of network evaluations --- the \ac{NFE} of the
sampler --- and in a receding-horizon controller that cost is paid again at every timestep.

This cost is not incidental but \emph{structural}, and it can be addressed where it arises. In a
diffusion model the number of steps is a
property of the trained noise schedule \parencite{ho2020denoising}; reducing it after training
degrades the sample, and recovering quality needs a different sampler
\parencite{song2021denoising}. In a deterministic transport model \parencite{lipman2023flow} the same
quantity is a property of the \ac{ODE} solver, and is chosen at inference. This thesis moves the
step budget from training time to inference time and measures what the move costs.

Once sampling is cheap, constraint enforcement becomes the dominant cost, and how constraints are
enforced becomes a question in its own right.

\section{Problem Statement}\label{sec:intro:problem}

The object of study is a closed-loop controller in which, at each timestep, a conditional
generative model proposes a trajectory over a fixed horizon, a constraint mechanism makes
that trajectory admissible, and the first action is executed. This thesis holds that
architecture fixed and varies two parts of it: the generative model and the projection method.

\emph{Flow matching} in the title names the family of generative models that transport noise to data
along an ordinary differential equation. The three generative models studied here belong to it and
differ in the velocity their network learns: \emph{instantaneous-velocity matching},
\emph{analytic average-velocity matching} and \emph{consistency-interpolated average-velocity
matching}. The study covers two observation modalities, state and camera images, and two embodiments, a manipulator
and a quadrotor.

The generative models and projection methods are first compared with the diffusion baseline on the
obstacle-avoidance benchmark of \ac{DPCC} \parencite{romer2025diffusion,jia2024towards}, with its
setup unchanged, and are then tested in a vision-conditioned alignment task and a quadrotor benchmark
built for this thesis. \autoref{ch:setup} describes the evaluation.

\section{Research Questions}\label{sec:intro:rqs}

\begin{description}
  \item[RQ1 --- Generative model.] With a matched temporal U-Net backbone, dataset and evaluation
    protocol, can the flow-based models reach the diffusion baseline's task success and constraint
    satisfaction with fewer sampling steps and less planning time per control step? How do the three
    objectives compare as the sampling budget changes?
  \item[RQ2 --- Constraints.] Given one fixed feasible set and one solver, does endpoint projection
    reach the goal with fewer constraint violations, fewer control steps or less computation than
    the per-step projection of \ac{DPCC}? For which step budgets and activation thresholds does it
    guide subsequent sampling, and when does it reduce to projection of the finished sample?
  \item[RQ3 --- Transfer.] Obstacle avoidance is planar, observed as a state, kinematic at the
    planning level, and solvable by a classical planner with a collision constraint. Do the answers
    to RQ1 and RQ2 hold when the observation becomes camera images, when the task requires aligning
    an object through contact, and when the platform is an underactuated, open-loop unstable
    quadrotor moving through three-dimensional obstacle fields --- settings whose dynamics and
    perception are closer to those of a physical deployment?
\end{description}

\srcnote{TARGET \S7 maps each RQ to its evidence sections; the former RQ2 on the step budget is
folded into RQ1 (v2.22).}

\section{Contributions}\label{sec:intro:contributions}

This thesis makes six contributions, in two stages. The first four are made on the obstacle-avoidance
benchmark of \ac{DPCC} \parencite{romer2025diffusion,jia2024towards}, with its setup unchanged, against
its diffusion model and its projection: three generative models take the diffusion model's place one
after another, and a second projection method is set beside the inherited one. The last two take what
that benchmark establishes --- the generative models and both projection methods, in the
configurations that hold there --- into two settings built for this thesis: a task observed through
cameras, and a plant that flies.

\paragraph{On the benchmark of \ac{DPCC}.}
\begin{enumerate}
  \item \textbf{Flow matching in place of diffusion.} The diffusion model of \ac{DPCC}
    \parencite{romer2025diffusion} is replaced by instantaneous-velocity matching, the flow-matching
    objective \parencite{lipman2023flow,lipman2024guide}, which learns a velocity field and generates a
    trajectory by integrating an ordinary differential equation. Backbone, projection, data and
    evaluation protocol are kept, and the number of network evaluations becomes a choice made at
    deployment (\autoref{sec:res:fewstep}).
  \item \textbf{Analytic average-velocity matching.} Analytic average-velocity matching
    \parencite{geng2025meanflow} learns the average velocity over an interval instead of the
    instantaneous velocity, with a regression target constructed analytically through a
    Jacobian--vector product, so that a single network evaluation can cover the whole integration. It
    replaces instantaneous-velocity matching within the same architecture (\autoref{sec:res:fewstep}).
  \item \textbf{Consistency-interpolated average-velocity matching.} Consistency-interpolated
    average-velocity matching \parencite{zhang2025alphaflow} also learns the average velocity, with a
    target interpolated by a consistency step ratio between the instantaneous velocity and a
    stop-gradient prediction of the network itself; the ratio is annealed during training. Both
    average-velocity objectives achieve state-of-the-art one- and two-step image generation
    \parencite{geng2025meanflow,zhang2025alphaflow}; this thesis evaluates both as planners inside a
    constrained controller (\autoref{sec:res:fewstep}).
  \item \textbf{Endpoint projection against DPCC's per-step projection.} The central contribution of
    DPCC is a model-based projection applied to each denoising step \parencite{romer2025diffusion}.
    Endpoint projection \parencite{li2025hardflow} instead projects the trajectory the sampler is
    predicted to reach and steers the sampler towards it. Both are implemented on the same feasible set
    and solver and compared, together with the step budget below which endpoint projection reduces to
    projection after sampling (\autoref{sec:res:constraints}).
\end{enumerate}

\paragraph{Beyond it.}
\begin{enumerate}
  \setcounter{enumi}{4}
  \item \textbf{A vision-conditioned alignment task.} Once established on the state-based benchmark
    of \ac{DPCC}, all generative models and both projection methods are tested on a more complex
    task built for this thesis from the alignment task of D3IL, in which the robot observes two camera
    views. The generative network is given a visual encoder, and the projection is given the dynamics
    rows and the constraint set of the new task (\autoref{sec:res:visual}).
  \item \textbf{A quadrotor benchmark.} The controller is carried to an underactuated, open-loop
    unstable plant: the manipulator with inverse kinematics is replaced by the Skydio~X2 quadrotor of
    MuJoCo Menagerie \parencite{zakka2022menagerie}, commanded through its rotor thrusts by a cascaded
    geometric tracking controller \parencite{lee2010geometric}, and the stack is tested from state
    observations in three scenes that ask three different questions. UAV-pillars is
    D3IL-avoiding flown by the quadrotor: its field is scaled into an arena of pillars, and
    the planner of the first four contributions flies there unchanged, so the scene reads what the
    vehicle adds to a plan D3IL-avoiding has already judged. UAV-corridor is the three-dimensional test:
    its two test-time constraints leave the horizontal plane, and a plan has to descend or climb to
    satisfy them. UAV-s-curve is the scene of two sharp turns, and it is reported for what those turns
    do to the tracking controller. This is where the constraint mechanism is asked to hold on a plant
    that requires continuous stabilisation: an arm can hold a setpoint through joint feedback, while
    a quadrotor must regulate both position and attitude in flight
    (\autoref{sec:bg:quadrotor}). The scenes, their constraint sets, the flown demonstrations of the
    corridor and the s-curve, and the control loop were built for this thesis in MuJoCo
    \parencite{todorov2012mujoco}, inside the environment code of D3IL; MuJoCo MPC
    \parencite{howell2022predictive}, a sampling-based predictive controller, is the second tracking
    controller, used in the controller comparison alone (\autoref{sec:res:uav}).
\end{enumerate}

\section{Outline}\label{sec:intro:outline}

\autoref{ch:background} introduces generative trajectory models, constraint enforcement and the two
control platforms. \autoref{ch:related} positions the comparison among generative planners and
constrained control methods. \autoref{ch:method} develops the models, projection procedures and
execution loops. \autoref{ch:setup} specifies the tasks, training and evaluation protocols.
\autoref{ch:results} reports the model and projection comparisons and their transfer to visual and
aerial tasks. \autoref{ch:conclusion} summarises the findings and answers the research questions, and
\autoref{ch:discussion} states the limits of the work, what a deployment would add to it, and future
work.

% =============================================================================
